\section{Introduction: Visual Decision Updates}

A VLM can answer an original image correctly and still fail to update that
decision after the visual evidence changes. CertVIC focuses on this narrower and
testable property: the visual decision update. The unit of evidence is not an
image, an edited image, or a model output in isolation. It is a certified
original/edit pair whose confounds have been checked before scoring.

The central measurement is the intervention-consistency gap: original-image
accuracy minus the rate at which the original and edited answers change in the
direction required by the intervention. This gap is only paper-relevant after
item-validity certification. Naive all-item gaps, smoke outputs, planned
artifacts, and unreviewed edits are analysis scaffolds, not evidence.

The tiny pilot is therefore decisive. Before any VLM inference begins, edit
detectability and visual quality must pass. If edit detectability is high, the
paper does not have a clean decision-update measurement; it has an artifact
confound. If detectability and item validity pass, CertVIC can proceed to
open-model inference and anytime-valid certification with claims scoped to the
models actually run.

Empirical findings remain [RESULT REQUIRED].
